\section*{Capítulo \ref{ch:g2:water-on-earth} --- El agua en la Tierra}

\begin{solution}{exo:g2:water-on-earth:1}
\omterm{def:g2:water-on-earth:salt-fresh}{Agua salada}: el \omterm{def:g2:water-on-earth:ocean}{mar}, el \omterm{def:g2:water-on-earth:ocean}{océano}. \omterm{def:g2:water-on-earth:salt-fresh}{Agua dulce}: el \omterm{def:g1:landscapes-around-us:water}{río}, la lluvia, el \omterm{def:g1:landscapes-around-us:water}{lago} de
\omterm{def:g1:landscapes-around-us:landform}{montaña}, el hielo de la \omterm{def:g1:landscapes-around-us:landform}{montaña}.
\end{solution}

\begin{solution}{exo:g2:water-on-earth:2}
Porque casi toda su superficie está cubierta por el agua azul de los
\omterm{def:g2:water-on-earth:ocean}{océanos}: vista desde el espacio, parece azul.
\end{solution}

\begin{solution}{exo:g2:water-on-earth:3}
$100 - 71 = 29$ cuadrados son tierra.
\end{solution}

\begin{solution}{exo:g2:water-on-earth:4}
Un \omterm{def:g2:water-on-earth:ocean}{océano} es más grande: un \omterm{def:g2:water-on-earth:ocean}{mar} es una extensión más pequeña de \omterm{def:g2:water-on-earth:salt-fresh}{agua
salada}, a menudo en el borde de un \omterm{def:g2:water-on-earth:ocean}{océano}.
\end{solution}

\begin{solution}{exo:g2:water-on-earth:5}
No: el agua del \omterm{def:g2:water-on-earth:ocean}{océano} es \omterm{def:g2:water-on-earth:salt-fresh}{agua salada}, y beberla hace que el cuerpo pierda
agua. El marinero tiene que llevar \omterm{def:g2:water-on-earth:salt-fresh}{agua dulce}.
\end{solution}

\begin{solution}{exo:g2:water-on-earth:6}
Por ejemplo: en los \omterm{def:g1:landscapes-around-us:water}{ríos}, en los \omterm{def:g1:landscapes-around-us:water}{lagos}, en el hielo y la nieve de las
\omterm{def:g1:landscapes-around-us:landform}{montañas} y de la Antártida, bajo el suelo, en la lluvia.
\end{solution}

\begin{solution}{exo:g2:water-on-earth:7}
\omterm{def:g2:water-on-earth:salt-fresh}{Agua salada}: 96 cubos. \omterm{def:g2:water-on-earth:salt-fresh}{Agua dulce}: $2 + 2 = 4$ cubos.
(En rigor, una parte del agua bajo el suelo también es salada; el dibujo la
cuenta con el \omterm{def:g2:water-on-earth:salt-fresh}{agua dulce}.)
\end{solution}

\begin{solution}{exo:g2:water-on-earth:8}
Los \omterm{def:g2:water-on-earth:ocean}{océanos} (96 cubos), después el hielo y la nieve (2 cubos), después los
\omterm{def:g1:landscapes-around-us:water}{ríos} (mucho menos de un cubo).
\end{solution}

\begin{solution}{exo:g2:water-on-earth:9}
Puede meterse en el suelo; puede correr por el suelo hasta un arroyo o un
\omterm{def:g1:landscapes-around-us:water}{río}; puede evaporarse y volver al aire cuando sale el sol.
\end{solution}

\begin{solution}{exo:g2:water-on-earth:10}
En el fondo del recipiente queda sal. El agua se ha evaporado: se ha ido al
aire en forma de vapor de agua.
\end{solution}

\begin{solution}{exo:g2:water-on-earth:11}
Por ejemplo: la gota baja por el tejado hasta el canalón y el desagüe,
después llega a un arroyo y a un \omterm{def:g1:landscapes-around-us:water}{río}; el \omterm{def:g1:landscapes-around-us:water}{río} la lleva al \omterm{def:g2:water-on-earth:ocean}{mar}; un día de sol
se evapora del \omterm{def:g2:water-on-earth:ocean}{mar}, sube al aire y pasa a formar parte de una nube.
\end{solution}
